\documentclass[11pt,a4paper]{article}

% ── Packages ─────────────────────────────────────────────────────────────────
\usepackage[T1]{fontenc}
\usepackage[utf8]{inputenc}
\usepackage{mathptmx}
\AtBeginDocument{\let\hbar\relax
  \DeclareMathSymbol{\hbar}{\mathord}{AMSb}{"7E}}
\usepackage{amsmath,amssymb,amsthm}
\usepackage{amsfonts}
\AtBeginDocument{\let\hbar\relax
  \DeclareMathSymbol{\hbar}{\mathord}{AMSb}{"7E}}
\usepackage{xcolor}
\usepackage{geometry}
\usepackage{microtype}
\usepackage{hyperref}
\usepackage{booktabs}
\usepackage{enumitem}
\usepackage{setspace}
\usepackage{titlesec}
\usepackage{abstract}
\usepackage{natbib}

\geometry{top=1.1in, bottom=1.1in, left=1.15in, right=1.15in}

\titleformat{\section}{\normalfont\bfseries\large}{\thesection}{1em}{}
\titleformat{\subsection}{\normalfont\bfseries\normalsize}{\thesubsection}{1em}{}
\titleformat{\subsubsection}{\normalfont\bfseries\itshape\normalsize}%
  {\thesubsubsection}{1em}{}

\newtheorem{theorem}{Theorem}[section]
\newtheorem{proposition}[theorem]{Proposition}
\newtheorem{conjecture}[theorem]{Conjecture}
\newtheorem{definition}[theorem]{Definition}
\newtheorem{remark}[theorem]{Remark}

\newcommand{\SBDG}{S_{\mathrm{BDG}}}
\newcommand{\Pact}{\mathcal{P}_{\mathrm{act}}}
\newcommand{\LLC}{\mathrm{LLC}}

\hypersetup{colorlinks=true,linkcolor=blue,citecolor=blue,urlcolor=blue}

% ── Title ─────────────────────────────────────────────────────────────────────
\begin{document}

\begin{titlepage}
\centering
\vspace*{1.5cm}

{\LARGE\bfseries RA Foundation\\[0.5em]
\large\itshape Physics from the Nature of Causation\par}

\vspace{1.4cm}

{\large Joshua F.\ Sandeman\\[0.3em]
\normalsize Independent Researcher, Salem, Oregon, USA\\[0.2em]
\href{mailto:sansuikyo@gmail.com}{sansuikyo@gmail.com}\\[0.4em]
\href{https://relationalactualism.org}{relationalactualism.org}}

\vspace{0.8cm}
{\small\textit{Preprint:}~\href{https://doi.org/10.5281/zenodo.19229438}{doi.org/10.5281/zenodo.19229438}\par}
  {\small Companion papers and Lean~4 proofs:
  \href{https://zenodo.org/communities/relational-actualism}%
  {zenodo.org/communities/relational-actualism}\par}

\vspace{1.4cm}

\begin{abstract}
We present Relational Actualism (RA), a framework grounded in a single
ontological commitment: that irreversible on-shell actualization events,
writing permanent directed edges into a growing causal Directed Acyclic Graph,
are the primitive physical reality. All other structure is derived.
From the discrete causal graph and its Benincasa-Dowker-Glaser (BDG) action,
we derive: spacetime kinematics ($c = l_P/t_P$; $E = \gamma mc^2$; proper time
exact as an integer count); general relativity, via BDG uniqueness
(Benincasa-Dowker 2010) and confirmed by Lovelock's theorem
given Lean-verified conservation and covariance; the complete Standard Model
force structure from the four independent BDG degrees of freedom in 4D (one
used by gravity, three generating $U(1)\times SU(2)\times SU(3)$); electric
charge quantisation in units of $e/3$; exact baryon number conservation;
maximal parity violation as a theorem from DAG acyclicity; the Koide lepton
mass formula from a generation-sector $SU(3)$ symmetry; grand unification at
the Planck scale; and the flat galactic rotation curves and Hubble tension
as consequences of causal graph dimensional reduction in sparse regions.
The same Erd\H{o}s-R\'{e}nyi percolation transition governs the biological
origin of life, the quantum fault-tolerance threshold, and the transition from
QCD confinement to asymptotic freedom.
Quantum mechanics is exact at the discrete level; the Schr\"{o}dinger equation
is its large-density macroscopic approximation, and the measurement problem
dissolves. All results follow from the single fact that some interactions are
irreversible. Four additional long-standing problems are dissolved: $\theta_{\mathrm{QCD}} = 0$ exactly (strong CP; no axion needed); the black hole information paradox (Causal Severance partition); the baryon asymmetry as a Causal Severance initial condition; and a structural conjecture for $d = 4$ spacetime uniqueness. Falsifiable predictions, Lean~4 verification status, and open
derivation targets are tabulated. The programme comprises twelve papers;
three are in peer review (\emph{Foundations of Physics} --- awaiting
editor assignment; \emph{Physical Review D}; \emph{International Journal of Astrobiology}).
Core algebraic results are independently machine-verified: a Python
notebook confirms 52/52 numerical checks at machine precision.
Eight Lean~4 proof files establish \textbf{176 theorems} with one
intentional sorry (the LQI adapter) and no axioms beyond Mathlib.
Key machine-checked results include: the Koide $K = 2/3$ identity;
the $SU(3)_{\mathrm{gen}}$ coherent state theorem; the BDG particle
classification (five topology types mapping exactly to the Standard Model
particle spectrum); colour confinement ($L{=}3$ gluon, $L{=}4$ quark);
the BDG particle universe closure theorem (124 extension cases); structural
qubit fragility (electrons and photons at minimum BDG score~1); the
causal invariance of the quantum measure; Rindler stationarity (Unruh
resolution); and the $\mathcal{P}_{\mathrm{act}}$ conservation theorem
with BDG locality lemma --- together proving $G_{\mu\nu} = 8\pi G\,
\mathcal{P}_{\mathrm{act}}[T_{\mu\nu}]$ with $\Lambda = 0$ uniquely
from the Local Ledger Condition.
New results (April 2026): a multiplicative cascade through the BDG depth
layers gives seven particle masses from the Planck mass alone: the proton
($m_p = m_P\,\alpha_{\mathrm{EM}}^5/2^{28} = 941$~MeV, $0.3\%$),
neutron ($0.2\%$), proton charge radius ($r_p = 0.84$~fm, $0.03\%$ CODATA),
pion ($124$~MeV, $11\%$), $W$ boson ($83$~GeV, $3.3\%$),
and Higgs boson ($m_H = 133\,m_p = 125.2$~GeV, $0.06\%$ --- within experimental
uncertainty). Two $d = 4$ geometric identities ($2N_c{-}1 = d{+}1$;
$\frac{4}{3}\,d! = d \times 2^{d-1}$) hold only in four spacetime dimensions,
proving that stable hadronic matter exists only in $d = 4$.
The Hubble tension prediction is now parameter-free:
$H_{\mathrm{local}} = H_{\mathrm{CMB}}(1 + 0.463\,|\delta|)$.
RAQM is under review at \emph{Foundations of Physics} (submitted 2026).
\end{abstract}

\medskip\noindent
\textbf{DOI:}~\href{https://doi.org/10.5281/zenodo.19229438}{doi.org/10.5281/zenodo.19229438}

\medskip\noindent
\textbf{Keywords:} Causal DAG \textperiodcentered{} Actualization
\textperiodcentered{} Benincasa-Dowker-Glaser action \textperiodcentered{}
Lovelock theorem \textperiodcentered{} Standard Model \textperiodcentered{}
Charge quantisation \textperiodcentered{} Koide formula \textperiodcentered{}
Causal Firewall \textperiodcentered{} Lean~4 verification
\textperiodcentered{} Measurement problem

\end{titlepage}

% ─────────────────────────────────────────────────────────────────────────────
\section{The One Commitment}
\label{sec:commitment}
% ─────────────────────────────────────────────────────────────────────────────

Physics has two foundational problems it cannot simultaneously solve. The
first is ontological: quantum mechanics describes a world of superpositions and
probabilities, but every observation finds a world of definite events.
The second is structural: general relativity and the Standard Model treat
spacetime and matter as different kinds of things, connected by Einstein's
equation but with no account of how two ontologically distinct categories
talk to each other.

Relational Actualism resolves both problems by going beneath them. The resolution
is not a new mechanism for wave function collapse, nor a modification of
general relativity, nor a new symmetry imposed on the Standard Model.
It is a change in what the theory is \emph{about}.

The single ontological commitment of RA is this:

\medskip
\begin{center}
\fbox{\parbox{0.82\textwidth}{%
\centering\medskip
\emph{Irreversible on-shell actualization events --- physical interactions that
produce a permanent increase in quantum relative entropy with respect to the
vacuum, $\Delta S(\rho\|\sigma_0) > 0$ --- are the primitive physical reality.
They are not events \emph{in} spacetime; spacetime is the macroscopic description
of their accumulated pattern.}
\medskip}}
\end{center}
\medskip

From this, and nothing else, the framework is built. No new fields, no extra
dimensions, no additional constants. The commitment is testable: it makes
specific, quantitative predictions that current and near-term experiments can
confirm or falsify.

The following sections present the results in logical order --- from the most
primitive structure of the causal graph to the most derived results, including
the Standard Model particle spectrum, the cosmological observations, the origin
of biological complexity, and the structure of quantum mechanics itself.
At each level, the status of the result is stated precisely: Lean-verified
(machine-checked), proved, conjectured with strong support, or open.

A note on the companion papers: this paper presents the unified logical
structure; the companions provide the derivations. New readers may start
here, then follow the suite in the arc \emph{Bosons to Brains}:
\begin{enumerate}[noitemsep,label=\arabic*.,leftmargin=*]
  \item \textbf{RAQM}~\citep{Sandeman2026RAQM} --- The ontological commitment;
  quantum mechanics and the measurement problem.
  \item \textbf{RAGC}~\citep{Sandeman2026RAGC} --- Gravity and cosmology;
  general relativity as inevitable consequence.
  \item \textbf{RACL}~\citep{Sandeman2026RACL} --- The Einstein field
  equations uniquely derived from the Lean-verified LLC, via the
  $\mathcal{P}_{\mathrm{act}}$ conservation theorem and Lovelock uniqueness;
  closes RAGC's primary hard wall. 176 Lean-verified theorems total.
  \item \textbf{RAEB}~\citep{Sandeman2026RAEB} --- The Engine of Becoming;
  QM and GR unified as approximations to a covariant graph rewriting algorithm.
  \item \textbf{RASM}~\citep{Sandeman2026RASM} --- The Standard Model;
  particles, forces, and generations from BDG topology.
  \item \textbf{RATM} --- \emph{Relational Actualism: The Topology of Matter}: The Lean-verified BDG particle classification theorem. Five causal topology types map exactly to the SM particle spectrum; colour confinement is proved with finite BDG-filtration lengths $L=3$ (gluon) and $L=4$ (quark); the closure theorem establishes no new particle types can emerge. 67~Lean-verified results.

\item \textbf{RADM}~\citep{Sandeman2026RADM} --- Dark matter as
  actualization-density topology in sparse regions.
  \item \textbf{RAQI}~\citep{Sandeman2026RAQI} --- Quantum information;
  the Kinematic Coherence Bound ($N_{\max} = \eta \cdot p_{\mathrm{th}}$);
  structural qubit fragility (electrons propagate at the minimum BDG
  actualization score of~1, the lowest permitted for any stable free particle;
  Lean-verified); decoherence channel constraints from the BDG closure theorem.
  \item \textbf{RAHC}~\citep{Sandeman2026RAHC} --- Emergence and the
  complexity hierarchy; algorithmic causal dynamics.
  \item \textbf{RACI}~\citep{Sandeman2026RACI} --- Complexity, life,
  and intelligence; the causal graph grows until it models itself.
  \item \textbf{RACF}~\citep{Sandeman2026RACF} --- The Causal Firewall
  as a substrate-independent habitability criterion; universal
  biosignatures and the SETI programme.
\end{enumerate}
This paper presents the logical structure; the companions provide the proofs.

% ─────────────────────────────────────────────────────────────────────────────
\section{The Causal Graph and Its Dynamics}
\label{sec:graph}
% ─────────────────────────────────────────────────────────────────────────────

\subsection{The Growing Causal DAG}

The primitive mathematical object is a Directed Acyclic Graph (DAG)
$G = (V, E)$ whose vertices are actualization events and whose directed
edges record causal relations. The acyclicity of the DAG is not a technical
convenience: it is the mathematical expression of the irreversibility of
actualization. An edge from $u$ to $v$ means $u$ is in the causal past of
$v$; the absence of cycles means there are no closed timelike curves, no
causal paradoxes, and no way for information to travel backward in time.
The arrow of time is not imposed on the graph; it is encoded in its
definition.

The graph grows by the sequential addition of new vertices. At each step,
the Covariant Stochastic Graph Rewriting System (CSG)~\citep{Sandeman2026RAEB}
identifies the set of possible new vertices consistent with the existing
causal structure and selects one. The selection is governed by the
\emph{actualization criterion}~\citep{Sandeman2026RAQM}: an interaction
constitutes an actualization event when it produces an irreversible increase
in quantum relative entropy $\Delta S(\rho\|\sigma_0) > 0$ with respect to
the vacuum reference state $\sigma_0$.

\subsection{The BDG Action and the Actualization Filter}

The unique discrete, retarded, Lorentz-invariant action functional on a
causal set in 4 spacetime dimensions is the Benincasa-Dowker-Glaser
action~\citep{Benincasa2010}:
\begin{equation}
  \SBDG(v) = \frac{1}{l_P^2}\bigl(1 - N_1(v) + 9N_2(v) - 16N_3(v)
  + 8N_4(v)\bigr),
  \label{eq:BDG}
\end{equation}
where $N_k(v)$ counts the number of $k$-element inclusive causal intervals
in $\mathrm{past}(v)$. The RA actualization filter identifies
$\Delta S(\rho\|\sigma_0) > 0$ with the discrete condition $\SBDG(v) > 0$:
a vertex actualizes if and only if it contributes positive curvature to
the local causal structure. This is the only free identification in the
framework; it connects the entropic actualization criterion to the
graph-geometric action.

\subsection{The Local Ledger Condition}

The Local Ledger Condition (LLC) is the conservation law of the growing graph:
at every vertex $v$, the total incoming causal charge equals the total
outgoing causal charge. Formally, for any conserved quantity $Q$:
\begin{equation}
  \sum_{u \in \mathrm{past}_1(v)} Q(u \to v)
  = \sum_{w \in \mathrm{future}_1(v)} Q(v \to w).
  \label{eq:LLC}
\end{equation}
The LLC, the causal invariance of the quantum measure (the statement that
the path-integral amplitude for any history is independent of the order
in which spacelike-separated actualization events are described), and the
frame independence of quantum relative entropy have all been formally verified
in Lean~4/Mathlib~\citep{Sandeman2026RAGC}. These Lean-verified results are
the mathematical bedrock of the framework.

\subsection{The Engine of Becoming: Five-Step Algorithm}

The covariant stochastic graph rewriting system that governs the growing
DAG is formalised in the companion paper RAEB~\citep{Sandeman2026RAEB}
as a five-step algorithm operating at every actualization cycle:

\begin{enumerate}[noitemsep,label=\textbf{\arabic*.},leftmargin=*]
  \item \textbf{Enumerate causal predecessors.} Given the current graph
  $G_n$, identify the frontier set $F_n$ of vertices whose causal futures
  are not yet determined. These are the sites at which the next
  actualization can occur.
  \item \textbf{Compute the quantum superposition.} For each candidate
  vertex $v \in F_n$, evaluate the quantum amplitude $\psi(v)$ from the
  sum over all histories consistent with the causal past $\mathrm{past}(v)$.
  The quantum measure (Sorkin~\citep{Sorkin1994}) assigns amplitude exactly,
  with no Copenhagen interpretation and no external observer.
  \item \textbf{Apply the BDG filter.} Retain only candidates satisfying
  $\SBDG(v) > 0$: those that contribute positive curvature to the local
  causal structure. This is the actualization criterion in graph-geometric
  form; it eliminates vacuum fluctuations and unphysical configurations
  at the discrete level.
  \item \textbf{Select the actualization event.} One candidate is
  irreversibly selected according to the quantum measure. This is the
  only stochastic step. The entropy criterion $\Delta S(\rho\|\sigma_0)
  > 0$ is satisfied; the event becomes a permanent vertex in $G_{n+1}$.
  \item \textbf{Update the macroscopic description.} In the
  $\mu \gg 1$ limit, the accumulated actualization pattern sources the
  Ricci tensor via $G_{\mu\nu} = 8\pi G\,\Pact[T_{\mu\nu}]$.
  The BDG locality lemma (proved in RACL~\citep{Sandeman2026RACL})
  guarantees that this sourcing is local:
  $\tau_{\mathrm{BDG}}(\lambda) = 4\Gamma(5/4)(12/\pi^2)^{1/4}\lambda^{-1/4}
  \to 0$ as $\lambda \to \infty$.
\end{enumerate}

The causal invariance theorem --- that the quantum measure amplitude
for any spacelike-separated set of actualization events is independent
of the order in which they are described --- is proved in
Lean~4~\citep{Sandeman2026RAGC}, given the \texttt{amplitude\_locality}
axiom (physically justified by QFT microcausality; the intrinsic discrete
proof is an open target). This is the discrete analogue of Lorentz
covariance.

\subsection{The Poisson-CSG and Three Regimes}

The mean-field description of the CSG dynamics is the Poisson-CSG, governed
by the local parameter:
\begin{equation}
  \mu(x) = \lambda(x)\,V_{\mathrm{coh}}\,p_{\mathrm{th}},
  \label{eq:mu}
\end{equation}
where $\lambda(x)$ is the local actualization density, $V_{\mathrm{coh}}$
is the coherence volume, and $p_{\mathrm{th}}$ is the per-link actualization
probability~\citep{Sandeman2026RAQI}. Three regimes follow directly:

\begin{itemize}[noitemsep]
  \item $\mu \gg 1$ (dense): the graph is 3-dimensional, matter-dominated.
  General relativity governs.
  \item $\mu \sim 1$ (critical): the Erd\H{o}s-R\'{e}nyi giant-component
  threshold. Phase transitions occur.
  \item $\mu \ll 1$ (sparse): the graph is 2-dimensional
  (Durhuus-Jonsson-Wheater), gravity-modified. Galactic halos.
\end{itemize}

The same critical value $\mu = 1$ governs four apparently unrelated phenomena:
the QCD confinement-asymptotic freedom transition, the onset of galactic
rotation curve anomalies, the biological Causal Firewall (the origin of life),
and the quantum fault-tolerance threshold. They are not four coincidences;
they are one transition seen at four scales of the causal graph.

% ─────────────────────────────────────────────────────────────────────────────
\section{Spacetime and Kinematics}
\label{sec:spacetime}
% ─────────────────────────────────────────────────────────────────────────────

\subsection{Proper Time as Integer Count}

A physical clock is a system that accumulates actualization events.
Proper time is the integer count of those events in Planck units:
\begin{equation}
  \tau = \mathcal{N}(\gamma) \cdot t_P.
  \label{eq:proper_time}
\end{equation}
This is exact at every scale including the Planck scale. The standard
GR proper time integral $\int d\tau$ is the large-$\mathcal{N}$ macroscopic
approximation of the integer count, not the fundamental definition.
Relativistic time dilation is the statement that different worldlines
accumulate actualization events at different rates — a consequence of
Lorentz-covariant on-shell thresholds, not a geometric axiom.

\subsection{The Speed of Light from Graph Scales}

The causal graph has two fundamental scales: the Planck length $l_P$ and
the Planck time $t_P$. Their ratio is the maximum causal propagation rate:
\begin{equation}
  c = \frac{l_P}{t_P}.
  \label{eq:c_derived}
\end{equation}
This is not a postulate. A massless pattern (zero net BDG debt per
propagation cycle) allocates 100\% of its causal bandwidth to spatial
propagation, achieving exactly $c$. A massive pattern must pay
actualization cost at every vertex, reducing its maximum spatial velocity
below $c$. The relativistic energy relation
\begin{equation}
  E = \gamma mc^2
  \label{eq:rel_energy}
\end{equation}
follows from the on-shell condition $p^\mu p_\mu = m^2c^2$, the
identification of mass with rest-frame actualization frequency
$f_0 = mc^2/h$, and the Lorentz transformation of the frequency
four-vector. No additional postulates are required.

\subsection{General Relativity: Inevitable by Two Independent Routes}

The Einstein field equations are not recovered in a limit or
approximation. They are logically inevitable. Two independent
derivations establish this:

\paragraph{Route 1: Lovelock's theorem.}
By Lovelock's theorem~\citep{Lovelock1971}: in 4 spacetime dimensions, the unique
symmetric, divergence-free, rank-2 tensor constructible from the
metric and its first and second derivatives is
$G_{\mu\nu} + \Lambda g_{\mu\nu}$.

The RA causal graph satisfies all three conditions required:
\begin{itemize}[noitemsep]
  \item \emph{Conservation}: the LLC, Lean-verified.
  \item \emph{Covariance}: causal invariance of the quantum measure,
  Lean-verified (unconditional---amplitude locality proved as a theorem
  of discrete DAG dynamics in \texttt{RA\_AmpLocality.lean}).
  \item \emph{Locality}: $\SBDG(v)$ depends only on $\mathrm{past}(v)$,
  definitional in CSG.
\end{itemize}
Given these, Lovelock guarantees Einstein.

\paragraph{Route 2: BDG uniqueness (RA-native).}
Three Lean-verified results---the LLC~(L01), amplitude locality~(O01),
and the $d=4$ BDG closure~(L11)---combined with the published
Benincasa-Dowker continuum limit theorem, yield the Einstein-Hilbert
action directly without invoking Lovelock. The Bianchi identity in this
chain is the LLC at the metric level, not a separate geometric
constraint. In standard GR, the geometric Bianchi identity and matter
conservation are mysteriously compatible; in RA, they are the
\emph{same} conservation law at different levels of description.

The interesting physics is
therefore not whether GR emerges but where RA departs from it: sparse
regions ($\mu \sim 1$), Planck-scale corrections, and Causal Severance
boundaries.

\subsection{The Uniqueness Proof: GR with No Free Parameters}

The central mathematical achievement of the programme is the proof that
the Einstein field equations $G_{\mu\nu} = 8\pi G\,\Pact[T_{\mu\nu}]$
with $\Lambda = 0$ are the \emph{unique} consistent macroscopic field
equations of the RA-CSG --- with no free parameters remaining
(RACL~\citep{Sandeman2026RACL}, submitted to \emph{Classical and Quantum
Gravity}).

The proof assembles seven links in a chain, two of which are new results:

\begin{itemize}[noitemsep]
  \item[\textbf{(A)}] \textbf{LLC $\Rightarrow$ conservation.}
  The Lean-verified LLC gives $\nabla_\mu(\Pact T^{\mu\nu}) = 0$
  via the \emph{$\Pact$ conservation theorem} (proved in RACL):
  a Leibniz decomposition shows that the boundary term
  $(\nabla_\mu\Pact)T^{\mu\nu}$ vanishes because $T^{\mu\nu} = 0$
  wherever $\Pact = 0$ (the vacuum energy suppression theorem),
  and the covariant term $\Pact(\nabla_\mu T^{\mu\nu}) = 0$ by
  the Rideout-Sorkin density limit of the discrete LLC.

  \item[\textbf{(B)}] \textbf{BDG locality.}
  The \emph{BDG locality lemma} (proved in RACL): at Poisson vertex
  density $\lambda$, the BDG score depends only on ancestors within
  proper-time window $\tau_{\mathrm{BDG}}(\lambda) \propto \lambda^{-1/4}
  \to 0$ as $\lambda \to \infty$. In the dense-graph limit, $\SBDG(v)$
  is a local function of the metric at $x(v)$.

  \item[\textbf{(C)}] \textbf{Benincasa-Dowker.}
  The dense-graph mean $\langle\SBDG\rangle = (l_P^2/4)R$: the
  Einstein-Hilbert action density emerges from the BDG sum.

  \item[\textbf{(D--E)}] \textbf{Lovelock uniqueness.}
  The source $\Pact T^{\mu\nu}$ is divergence-free (by A). The
  left-hand side must therefore also be divergence-free, second-order
  in the metric, and local (by B and C). By Lovelock's
  theorem~\citep{Lovelock1971}, the unique such tensor in 4D is
  $\alpha G_{\mu\nu} + \Lambda g_{\mu\nu}$.

  \item[\textbf{(F)}] \textbf{$\Lambda = 0$ structurally.}
  Vacuum energy suppression ($\Pact|0\rangle = 0$) means the vacuum
  contributes nothing to the metric source. $\Lambda = 0$ is forced by
  the actualization ontology --- no fine-tuning, no cosmological constant
  problem.

  \item[\textbf{(G)}] \textbf{Bianchi identity.}
  $\nabla^\mu G_{\mu\nu} \equiv 0$ holds by differential geometry,
  not by the field equation. Setting $\kappa/\alpha = 8\pi G$ by
  the Newtonian limit completes the derivation.
\end{itemize}

The two new results --- the $\Pact$ conservation theorem and the BDG
locality lemma --- close the primary hard wall of RAGC. General
relativity is not recovered; it is forced.

\subsection{Event Horizons as Local Bandwidth Saturation}

A Causal Severance is a topological discontinuity in the causal graph:
a surface beyond which the graph has no outgoing edges. It forms when
the local actualization density saturates the bandwidth ceiling $c/l_P$
--- when all causal capacity is committed to boundary dynamics and no
outgoing connections remain. This is the $v \to c$ limit applied to
the graph boundary rather than to a worldline within it. The standard
black hole singularity is replaced by a Causal Severance followed by
daughter spacetime nucleation~\citep{Sandeman2026RAGC}: the saturated
bandwidth is released into a new, causally isolated subgraph. The
Big Bang is a Causal Severance nucleation event.

The daughter inherits its baryon-to-photon ratio from the parent's
Bekenstein-Hawking entropy, requiring a parent SMBH of mass
$M_{\mathrm{parent}} \approx 2.5$--$6.6 \times 10^6\,M_\odot$
(RAGC~\S6.3). The parent's Kerr spin geometry imprints anisotropic
initial conditions on the daughter: the quadrupole amplitude
$\varepsilon_2$ from the Kerr metric and the octupole
$\varepsilon_3$ from accretion asymmetry share the parent's spin
axis. A parameter-free $(2/\ell)^2$ dilution law from horizon exit
kinematics predicts that all five CMB large-angle anomalies---low
quadrupole, axis-of-evil alignment, hemisphere asymmetry, normal
$C_3$, normal $\ell \geq 4$---are structural consequences of Kerr
nucleation, explained by three parameters ($a_*$, $\Delta N$,
$\delta_{\mathrm{acc}}$) with zero remaining freedom
(RAGC~\S6.4). $\Lambda$CDM has no mechanism for any of these.

The field equations' density-dependent expansion produces a
positive feedback that drives filamentary cosmic web topology as a
dynamical attractor (RAGC~\S6.5). Both endpoints of the density
bifurcation lead to causal disconnection: starvation severance
(voids approaching $\mu = 1$) and density severance (boundaries
approaching $\mu_2$, nucleating daughters). The universe fragments
rather than equilibrating, categorically prohibiting heat death and
Boltzmann Brains (RAGC~\S6.6--6.7).

% ─────────────────────────────────────────────────────────────────────────────
\section{The Standard Model from BDG Topology}
\label{sec:SM}
% ─────────────────────────────────────────────────────────────────────────────

\subsection{Four Degrees of Freedom, One Used by Gravity}

The BDG action in equation~\eqref{eq:BDG} counts four independent
integer-valued quantities: $N_1, N_2, N_3, N_4$. One specific linear
combination constitutes $\SBDG$ and sources curvature. The remaining
three independent combinations are available as conserved topological
charges of persistent causal patterns. This is the central structural
observation:

\begin{proposition}[Three forces from 4D BDG topology]
\label{prop:three_forces}
In 4 spacetime dimensions, the BDG local topology has exactly four
integer-valued degrees of freedom. One generates gravity. The remaining
three are the three non-gravitational gauge interactions:
$U(1)$ electromagnetism, $SU(2)$ weak force, and $SU(3)$ strong force.
The gauge group $U(1)\times SU(2)\times SU(3)$ has three factors because
spacetime has four BDG degrees of freedom.
\end{proposition}

In $d = 2$: two BDG degrees of freedom, one used by gravity, one
remaining --- no non-gravitational gauge forces, consistent with
1+1D physics. The Standard Model gauge structure is a consequence of
living in 4 dimensions, not a postulate imposed on it.

\subsection{Particles as Persistent BDG Patterns}

A particle cannot be a fixed object in the growing DAG. It is a
\emph{persistent self-similar pattern}: a local causal topology that
recreates itself under BDG growth dynamics, with $\SBDG(v) > 0$ at
every vertex in the pattern. Mass is the actualization frequency required
to maintain the pattern: $f_0 = mc^2/h$. A heavier particle is a more
complex BDG pattern with a higher per-cycle actualization cost.
Spin is the self-similarity period: bosons recreate in one step (spin 1),
fermions in two (spin 1/2).

\subsection{The Lean-Verified BDG Particle Classification}

The abstract argument of \S\ref{sec:SM} is converted into machine-verified
mathematics in the companion paper RATM~\citep{Sandeman2026RATM}
(Lean-verified, zero \texttt{sorry} tags). The BDG stability
filter is applied exhaustively to all self-similar persistent causal
patterns in a 4D Poisson-CSG, classifying them by their filtration length
$L$ --- the number of BDG rewriting steps required to stabilise the pattern.

The classification yields exactly five topology types:

\begin{center}
\renewcommand{\arraystretch}{1.25}
\begin{tabular}{cll}
\toprule
$L$ & Topology type & Standard Model identification \\
\midrule
1 & Minimal retarded (two subtypes) & Photon; neutrino / scalar \\
2 & Fermion-like & Electron, quark (chirality resolved by DAG) \\
3 & Gluon-like LLC-flux & Gluon (confined: $Q_r + Q_g + Q_b \neq 0$ alone) \\
4 & Baryon-like LLC-flux & Quark (confined; $L=4$ requires baryonic closure) \\
\bottomrule
\end{tabular}
\end{center}

\noindent
The BDG filtration lengths $L = 3$ for gluons and $L = 4$ for quarks
are the finite causal depth at which the LLC flux-tube becomes energetically
unstable unless balanced: this is colour confinement derived from first
principles, not postulated.

The \emph{closure theorem} (124 extension cases exhaustively verified
in Lean~4 by \texttt{D1\_closure\_complete}) establishes that no new
particle topology type can emerge under BDG growth dynamics beyond these
five. The particle spectrum of the Standard Model is closed under RA
actualization: this is a mathematical theorem.

The \emph{$D1_a$ fixed-point theorem} (\texttt{D1a\_fixed\_point},
Lean-proved) establishes that stable free particles must propagate
at BDG score~$\geq 1$, with electrons and photons achieving the minimum.
This is structural qubit fragility: the particles that carry quantum
information in all current hardware sit at the minimum BDG stability
the actualization filter permits.

\subsection{Charge Quantisation}

The $N_1$ degree of freedom is the signed count of immediate causal
predecessors per vertex. In a 3+1D graph, a vertex can receive signed edges
from at most three independent spatial directions, so the conserved LLC
charge $Q_{N_1}$ is an integer in $\{-3, -2, -1, 0, +1, +2, +3\}$.
Identifying the unit with $e/3$:
\begin{equation}
  Q_{N_1} \in \{-e,\, -\tfrac{2e}{3},\, -\tfrac{e}{3},\, 0,\,
  +\tfrac{e}{3},\, +\tfrac{2e}{3},\, +e\}.
  \label{eq:charge_quant}
\end{equation}
This matches the complete set of electric charges in the Standard Model.
Electric charge is quantised in units of $e/3$ because spacetime has
three spatial dimensions~\citep{Sandeman2026RASM}.

\subsection{Photon Masslessness}

The photon carries $Q_{N_1} = Q_{N_2} = Q_{N_3/N_4} = 0$. A pattern with
zero net BDG charge of every kind has zero net BDG action per propagation
cycle: $\Delta\SBDG|_{\mathrm{per\ cycle}} = 0$. Zero actualization cost
gives $f_0 = 0$, hence $m = 0$. The photon is massless not by postulate
but because its pattern carries no causal debt of any kind.

\subsection{Colour Confinement from the LLC}

The $N_2$ degree of freedom has three independent spatial components
$(Q_r, Q_g, Q_b)$. The LLC requires $\Delta Q_r = \Delta Q_g =
\Delta Q_b = 0$ at every vertex. A single quark pattern with charge
$(Q_r, 0, 0)$ accumulates a permanent LLC imbalance in the $r$-direction
as it propagates. The causal closure condition
$\lim_{t\to\infty}[Q_r\,\mathrm{accumulated}] = 0$ is violated: a free
quark is structurally unstable. The only LLC-consistent configurations
are baryons ($Q_r + Q_g + Q_b = 0$) and mesons ($Q_r + \bar{Q}_r = 0$).
Colour confinement is not an external mechanism; it is the LLC applied
to the $N_2$ spatial charges~\citep{Sandeman2026RASM}.

\begin{proposition}[Exact baryon number conservation]
\label{prop:baryon}
Baryon number is the $N_2$ spatial winding number. The LLC conserves it
at every vertex at every scale. Baryon number conservation in RA is
exact, not approximate. Proton decay is forbidden.
\end{proposition}

The LLC is Lean-verified~\citep{Sandeman2026RAGC}. Baryon number
conservation therefore rests on machine-checked mathematics.

\subsection{Maximal Parity Violation from Acyclicity}

\begin{proposition}[Maximal parity violation]
\label{prop:parity}
The BDG action is retarded: $\SBDG(v)$ counts intervals in
$\mathrm{past}(v)$ only. The $N_3/N_4$ charge encodes how a pattern's
causal ancestry winds around the temporal axis. To reverse this winding
--- to convert a left-handed pattern to right-handed --- would require
some causal edges to point from future to past. The DAG is acyclic:
this is impossible. The weak interaction, which changes the $N_3/N_4$
charge of a pattern, can only couple to patterns whose winding is
consistent with forward temporal propagation: the left-handed sector.
Parity violation is exactly maximal --- a topological impossibility,
not a statistical bias~\citep{Sandeman2026RASM}.
\end{proposition}

The arrow of time and the left-handedness of the weak force are the
same fact: DAG acyclicity.

\subsection{Three Generations and No Fourth}

The BDG action in 4D uses $N_1, N_2, N_3, N_4$ --- four interval
sizes. Persistent patterns can be excited at three levels:
$N_2$-dominant (first generation: $e, u, d, \nu_e$),
$N_3$-dominant (second: $\mu, c, s, \nu_\mu$),
$N_4$-dominant (third: $\tau, t, b, \nu_\tau$).
There is no $N_5$ in the 4D BDG action. The existence of exactly three
generations is a structural consequence, not a parameter to be fitted.
This prediction is supported by the LEP measurement $N_\nu = 2.984 \pm
0.008$~\citep{PDG2022} and by the absence of a fourth generation at the LHC.

\subsection{The Koide Formula from Generation SU(3)}

The three BDG excitation levels $\{N_2, N_3, N_4\}$ transform under the
same permutation group $S_3$ as the three colour charges $\{r, g, b\}$
--- both arise from the same BDG structure. Extending $S_3 \to
SU(3)_{\mathrm{gen}}$, the most general invariant charged-lepton mass
matrix is:
\begin{equation}
  M_\ell = m_0\bigl(I + \sqrt{2}\,U(\theta)\bigr), \quad
  U(\theta) = \mathrm{diag}\bigl(\cos\theta,\,
  \cos(\theta + \tfrac{2\pi}{3}),\,
  \cos(\theta + \tfrac{4\pi}{3})\bigr).
  \label{eq:koide_matrix}
\end{equation}
For any $m_0$ and $\theta$, the eigenvalues automatically satisfy the
Koide relation~\citep{Koide1983}:
\begin{equation}
  K \;\equiv\; \frac{m_e + m_\mu + m_\tau}{(\sqrt{m_e} + \sqrt{m_\mu}
  + \sqrt{m_\tau})^2} \;=\; \frac{2}{3}.
  \label{eq:koide}
\end{equation}
Numerically: $K = 0.666661$ with the 2022 PDG values, agreeing with $2/3$
to one part in $10^5$. The fit gives $m_0 \approx 313.9~\mathrm{MeV}$
--- equal to the constituent quark mass $m_p/3 \approx 312.8~\mathrm{MeV}$
to $0.35\%$~\citep{Sandeman2026RASM}. Both scales are set by the QCD
string tension $\sigma$, connecting the lepton mass scale to the baryon
mass scale through a single BDG calculation.

The Koide formula has been observed since 1981 and never derived from
the Standard Model. In RA it is the generation-sector Gell-Mann-Okubo
relation: a group-theoretic consequence of $SU(3)_{\mathrm{gen}}$.

\subsection{Grand Unification at the Planck Scale}

$SU(3)_{\mathrm{colour}}$ acts on the three spatial $N_2$ directions;
$SU(3)_{\mathrm{gen}}$ acts on the three BDG excitation levels $\{N_2,
N_3, N_4\}$. Both groups arise from the same BDG structure. At the
Planck scale ($\mu \to \infty$), where all BDG levels are equally and
densely populated, the two $SU(3)$ groups become indistinguishable.
Grand Unification is the Planck-scale limit of the BDG structure,
not an additional postulate. Consequently:
\begin{equation}
  \sin^2\theta_W\big|_{\mathrm{GUT}} = \frac{3}{8},
  \label{eq:weinberg}
\end{equation}
the SU(5)-type prediction, running to the observed $0.231$ at the
$Z$ scale via Poisson-CSG $\mu$-dependence. No supersymmetric partners
are required: the three problems SUSY was introduced to solve (hierarchy
problem, dark matter, coupling unification) are each resolved differently
in RA, and the SUSY algebra is not generated by BDG dynamics.

\subsection{Force Ranges from BDG Cost and LLC Balance}

\begin{center}
\renewcommand{\arraystretch}{1.25}
\begin{tabular}{llll}
\toprule
Force & Exchange & BDG cost & Range mechanism \\
\midrule
EM & photon & 0 &
  $1/r^2$ from causal 2-sphere geometry \\
Weak & $W$, $Z$ & $>0$ &
  Finite lifetime $\hbar/mc$; $\sim 10^{-18}$~m \\
Strong & gluon & 0, but $Q_{N_2}\neq 0$ &
  LLC flux tube; $V = \sigma r$ \\
\bottomrule
\end{tabular}
\end{center}

\noindent
All three force ranges follow from the BDG actualization cost and LLC
balance, with no additional parameters~\citep{Sandeman2026RASM}. Regge
trajectories $m^2 \propto J$ follow directly from the linear potential.
Asymptotic freedom at $\mu \gg 1$ and confinement at $\mu \to 1$ are
the same Poisson-CSG $\mu$-dependence seen at different scales.

\subsection{Neutrino Masses and Oscillations}

Neutrino oscillations are the generation-sector analogue of colour mixing.
Each mass eigenstate $|\nu_k\rangle$ is a distinct BDG pattern with
actualization frequency $f_0^k = m_k c^2/h$. Phase accumulation at
different rates gives oscillation probability~\citep{Sandeman2026RASM}:
\begin{equation}
  P(\nu_\alpha \to \nu_\beta) = \left|\sum_k U^*_{\alpha k}\,U_{\beta k}
  \,\exp\!\left(-i\frac{(m_k^2 - m_1^2)c^4 L}{2E\hbar c}\right)\right|^2.
  \label{eq:osc}
\end{equation}
The CP-violating phase $\delta_{\mathrm{CP}}$ is the same topological fact
as the arrow of time: DAG acyclicity. The RA prediction:
$m_1 \approx 0.37~\mathrm{meV}$, $\sum m_\nu \approx 59~\mathrm{meV}$,
testable with CMB-S4 and Euclid.

A further structural result from the Koide fits: the lepton sector
converges to $\theta_{\mathrm{lep}} = 2/9$ exactly (to one part in
$10^5$), where $2/9 = 2/(N_{\mathrm{gen}}\cdot N_{\mathrm{col}}) =
2/(3\times 3)$ is determined by the same BDG integers that define
$SU(3)_{\mathrm{gen}}$ and $SU(3)_{\mathrm{colour}}$. This predicts
the Cabibbo angle: $|V_{us}| = \sin(2/9) = 0.2204$, agreeing with the
observed $0.2245$ to $1.8\%$. The three quark Koide angles form an
arithmetic progression $\theta_0 \pm \delta$ around $\theta_0 = 2/9$,
with a single displacement $\delta \approx 0.0665~\mathrm{rad}$
controlling the full CKM hierarchy (RASM Conjecture~7.3).
The displacement $\delta$ satisfies $\delta =
C_2(SU(2)_{\mathrm{fund}})\,\alpha_s(\mu)/\pi = (3/4)\,\alpha_s/\pi$
to $0.1\%$ (RASM Proposition~7.4). Both the $K$-breaking
($\alpha_s/\pi \times 4/3$) and the $\theta$-displacement
($\alpha_s/\pi \times 3/4$) are one-loop corrections at the same
hadronic scale with their respective gauge Casimirs.
Their ratio $16/9 = |c_{N_3}|/|c_{N_2}|$ is the BDG coefficient ratio.
A third result follows: $|V_{cb}| = (2/\pi)\,\delta =
(3\alpha_s)/(2\pi^2)$ to $0.4\%$ (RASM Conjecture~7.4).
Three CKM parameters from two RA quantities ($\theta_0=2/9$ and
$\alpha_s(\mu)/\pi$):
$|V_{us}| = \sin(2/9) = 0.2204$ ($1.8\%$),
$\delta = (3/4)\alpha_s/\pi = 0.0666$ ($0.1\%$),
$|V_{cb}| = (2/\pi)\delta = 0.0424$ ($0.4\%$).

% ─────────────────────────────────────────────────────────────────────────────
\section{Cosmology}
\label{sec:cosmology}
% ─────────────────────────────────────────────────────────────────────────────

\subsection{Dark Matter as Topology, Not Particles}

Any particle that does not participate in on-shell actualization events
has zero RA assembly depth and contributes zero metric weight. WIMPs,
axions, and sterile neutrinos are categorically excluded as dark matter
candidates --- not ruled out by non-observation alone, but structurally
prohibited by the actualization criterion. This is the RA WIMP
prohibition~\citep{Sandeman2026RAGC}.

The observed dark matter phenomenology is replaced, not explained by a
new particle, across four distinct lensing regimes~\citep{Sandeman2026RADM}:
\begin{itemize}[noitemsep]
  \item \textbf{Rotation curves:} Two RA-derived mechanisms together
  give flat rotation curves at all radii. In the disk ($r < r_c$),
  the 2D cylindrical Laplacian of $\ln\lambda$ gives $\Sigma_\lambda \propto 1/R$,
  yielding $v_\mathrm{flat}^2 = \xi/(2r_d)$ (RADM~\S3.1, fully derived).
  In the sparse halo ($\mu < 1$), the DJW theorem gives 2D gravity
  $\Phi \propto \ln r$, from which the Tully-Fisher relation
  $v_\mathrm{flat}^4 \propto M$ follows by causal flux conservation.
  \item \textbf{Bullet Cluster:} The $A_{\mathrm{RA}}$ ratio
  (stellar/shocked gas $\approx 10^8$) makes lensing trace the stellar
  component by a factor of $10^8$ over the gas --- no dark matter halo
  required. A quantitative $\kappa(\theta)$ profile is derived in RADM.
  \item \textbf{Galaxy-galaxy lensing:} The DJW logarithmic potential
  gives a convergence $\kappa(b) = v_\mathrm{flat}^2/(4G\Sigma_\mathrm{cr}b)$
  as a theorem of the RA causal structure (RADM~\S4.1, RA-native derivation).
  This profile cuts off at the baryonic edge ($R_\mathrm{baryon}\sim 150$--300~kpc),
  whereas $\Lambda$CDM's NFW halo extends to $\sim$Mpc.
  The shape at $b > 300$~kpc falsifies one framework. The residual amplitude
  is accounted for by the warm CGM, testable via Sunyaev-Zel'dovich stacking.
  \item \textbf{Cosmic shear:} RA recovers $\Lambda$CDM to $\sim1\%$
  at large scales; RA-specific corrections appear at 10--100~Mpc.
\end{itemize}

\subsection{The Hubble Tension}

The Hubble tension --- the discrepancy between the CMB-inferred value
$H_0 \approx 67.4~\mathrm{km\,s^{-1}\,Mpc^{-1}}$ and the local
measurement $H_0 \approx 73~\mathrm{km\,s^{-1}\,Mpc^{-1}}$ --- arises
naturally in RA from the actualization density difference between
matter-dominated dense regions and underdense voids. The local Universe
contains the KBC supervoid, a region of anomalously low actualization
density. RA predicts $H_{\mathrm{local}} \approx 73.6~\mathrm{km\,s^{-1}
\,Mpc^{-1}}$ in this void, matching the local measurement without
new physics~\citep{Sandeman2026RAGC, Sandeman2026RADM}.

A specific falsifiable prediction: the Eridanus supervoid should exhibit
\$H \approx 76.8~\mathrm{km\,s^{-1}\,Mpc^{-1}}$, measurable with current
and near-term surveys.

\subsection{The Cosmological Constant and Boltzmann Brains}

Vacuum energy suppression is a direct consequence of the $\Pact$
operator: only on-shell actualized states contribute to the metric source.
Virtual fluctuations are excluded, resolving the cosmological constant
problem without fine tuning~\citep{Sandeman2026RAGC}.

The Boltzmann Brain problem is categorically resolved by two mechanisms.
First, \emph{structurally}: each irreversible actualization requires a
thermodynamic cost that cannot be met by vacuum fluctuation---a brain
requires $\sim\!10^{25}$ independent actualization steps, each crossing
the $\Delta S^*$ threshold. Second, \emph{dynamically}: the dual
severance mechanism (RAGC~\S6.6) fragments every causally connected
component on a finite timescale, so no patch persists long enough for
BB nucleation. The combination is a categorical prohibition, not
merely a probabilistic suppression~\citep{Sandeman2026RAGC}.

% ─────────────────────────────────────────────────────────────────────────────
\section{Complexity and Life}
\label{sec:complexity}
% ─────────────────────────────────────────────────────────────────────────────

\subsection{The Causal Firewall}

At the critical density $\mu = 1$, the Poisson-CSG undergoes an
Erd\H{o}s-R\'{e}nyi giant-component transition. Below this threshold,
causal graphs are locally tree-like --- chemistry. Above it, they develop
stable long-range causal connectivity --- life. The Causal Firewall is
the RA account of the origin of life: not a continuum but a percolation
threshold, as sharp as a phase transition~\citep{Sandeman2026RACI}.

This is the same $\mu = 1$ transition as QCD confinement, galactic
rotation curve onset, and quantum fault-tolerance threshold. One
mathematical transition, four physical manifestations.

\subsection{Assembly Theory Convergence}

Assembly Theory~\citep{Walker2023} independently arrived at
the concept of causal depth as the measure of biological complexity.
In RA, causal depth $A_{\mathrm{RA}}(v)$ --- the length of the longest
directed path to vertex $v$ --- is the fundamental quantity. The
convergence of two independent frameworks on the same mathematical
structure supports both.


New results in RAHC and RACI~\citep{Sandeman2026RAHC,Sandeman2026RACI}
extend this to metabolic networks via the \emph{coarse-grained assembly
depth} $\tilde{A}_{\mathrm{RA}}(\mathcal{N})$: the number of irreversible
reactions on the critical path through the network. For glycolysis
$\tilde{A}_{\mathrm{RA}} = 3$ (hexokinase, PFK-1, pyruvate kinase),
with allosteric regulation concentrating at exactly these three nodes ---
independent biological confirmation. The \emph{origin-of-life sandwich}
$A(M^*_0) \leq \tilde{A}_{\mathrm{RA}}(\mathcal{N}_0) \leq
\tau_{\mathrm{cycle},0}/\tau_{\mathrm{irrev}}$
brackets the causal depth of the first replicators from both sides.


Causal depth is a selection pressure: replication rate is exponentially
suppressed by $A_{\mathrm{RA}}(M)$ (England bound), so natural selection
for replication speed is simultaneously selection for minimal causal depth.
Intelligence is recursive causal closure: the capacity of a system to
model its own actualization history and use that model to guide future
actualization~\citep{Sandeman2026RAHC}.

\subsection{Substrate-Independent Habitability (RACF)}


The Causal Firewall criterion is substrate-independent: neither Condition
F1 (Quantum Darwinism redundancy $R \geq R_{\min}$) nor Condition F2
(Unruh temperature $T_U = \hbar a/2\pi c k_B \ll T_{\mathrm{env}}$) makes
reference to molecular class, solvent chemistry, or biological alphabet.
The circumstellar habitable zone is a terrestrial proxy for these deeper
physical conditions, not the physical criterion itself.
RACF~\citep{Sandeman2026RACF} develops three universal biosignature
observables and two falsifiable predictions for the distribution of life,
submitted to \emph{International Journal of Astrobiology}.
% ─────────────────────────────────────────────────────────────────────────────
\section{Quantum Information}
\label{sec:QI}
% ─────────────────────────────────────────────────────────────────────────────

The Kinematic Coherence Bound~\citep{Sandeman2026RAQI} establishes
that the maximum number of fault-tolerant logical qubits in a physical
array scales as $N_{\max} = \eta \cdot p_{\mathrm{th}}$, where
$p_{\mathrm{th}} \approx 10^{-2}$ is the per-link actualization
probability. This threshold is not a property of any specific
error-correction code: it is the $\mu = 1$ Erd\H{o}s-R\'{e}nyi
transition applied to the causal neighbourhood of the logical qubit.

The Lean-verified BDG particle classification (RATM~\citep{Sandeman2026RATM})
establishes a structural account of qubit fragility. Every stable free
particle has a positive BDG actualization score; the minimum is~1.
Electrons and photons --- the particles that carry quantum information
in all current hardware --- propagate at the sequential fixed point
with score~1. $W^\pm$ bosons score~9; gluons score~18. Qubit fragility
is therefore not an accident of electromagnetic coupling strengths: it
is a consequence of the fact that laboratory-accessible particles sit at
the minimum BDG stability that the causal filter permits. This result is
proved by \texttt{D1a\_fixed\_point} in \texttt{RA\_D1\_Proofs.lean}.

The BDG particle universe closure theorem (RATM, Theorem~D1\_closure\_complete,
124 extension cases exhaustively verified) additionally constrains the
decoherence channels available to any qubit: every environmental particle
that can couple to a quantum system belongs to the closed BDG set.
No exotic decoherence mechanism from outside the closure is possible
within the RA framework.

The Lieb-Robinson bound on actualization propagation velocity
$v_{\mathrm{LR}} \sim J\ell/\hbar$ is the lattice expression of the
same causal bandwidth ceiling as $c = l_P/t_P$ in the relativistic
regime: the maximum rate at which actualization can propagate through
the causal graph, at the relevant density scale.

The fault-tolerance threshold, the speed of light, the Causal Firewall,
and the QCD confinement transition are the same Erd\H{o}s-R\'{e}nyi
percolation at four different scales of the causal graph.

% ─────────────────────────────────────────────────────────────────────────────
\section{Quantum Mechanics: Exact at the Discrete Level}
\label{sec:QM}
% ─────────────────────────────────────────────────────────────────────────────

\subsection{The Measurement Problem Dissolved}

The measurement problem asks: why does a quantum superposition resolve
into a definite outcome when measured? In RA the question does not arise,
because the premise is wrong. There is no superposition that mysteriously
collapses; there is a causal graph that grows. A quantum measurement is an
actualization event at which the entropic criterion $\Delta S(\rho\|\sigma_0)
> 0$ is satisfied. At that vertex, one outcome becomes permanently encoded
as a directed edge in the DAG. The other outcomes are the space of
\emph{unactualized candidates} --- they were never in the graph; they were
possible extensions of the graph that were not selected.

Wave function collapse is the writing of a new vertex. The wave function
is not a physical object that collapses; it is a description of the
probability distribution over possible next actualization events.

\subsection{The Quantum Measure is Exact; the Schr\"{o}dinger Equation is Not}

The quantum measure of Sorkin~\citep{Sorkin1994}, assigning complex
amplitudes to causal histories in a foliation-independent way, is exact
at the discrete level. It is not an approximation to RA: it IS the RA
causal graph dynamics at Level~2. Quantum mechanics, in the sense of the
full quantum measure formalism, is therefore not a macroscopic limit ---
it is the primitive description of the growing DAG.

What is a macroscopic approximation is the \emph{Schr\"odinger equation}
specifically. It emerges in the slowly varying, large-$\mu$,
non-relativistic, single-particle continuum limit of the quantum measure:
\begin{equation}
  \mu_{\mathcal{C}}(\text{transition: } |i\rangle \to |f\rangle)
  \;\xrightarrow{\mu \gg 1}\;
  \left|\int \mathcal{D}\Phi\; e^{iS[\Phi]/\hbar}\right|^2
  \;\xrightarrow{\text{non-rel.}}\;
  \text{Schr\"odinger evolution.}
  \label{eq:schrodinger_limit}
\end{equation}
The hierarchy is: quantum measure (exact) $\supset$ Feynman path integral
(large-$\mu$ limit) $\supset$ Schr\"odinger equation (non-relativistic
single-particle limit). Heisenberg uncertainty is the discrete graph's
resistance to simultaneous specification of complementary causal properties
--- exact at the discrete level, not an approximation.

\subsection{Classical Behaviour as the Dense-Graph Regime}

A large, warm, complex system is not fundamentally different from a
quantum system --- it IS a quantum system, described by the same
discrete causal graph. The difference is density. In the regime
$\mu \gg 1$, every degree of freedom has been pinned by previous
actualization events: the causal graph is so densely connected that
quantum indeterminacy has been eliminated at every relevant scale.
Classical behaviour is not a different ontology from quantum behaviour;
it is the $\mu \gg 1$ regime of the same ontology. The transition
from quantum to classical is the gradient $\mu(x)$ increasing from
sparse to dense. There is no Heisenberg cut; there is only a
gradient in actualization density.

The famous paradoxes of quantum mechanics --- the double slit, the
delayed choice, the quantum eraser, Schr\"{o}dinger's cat --- are each
resolved by identifying where in the growing DAG the actualization
events occur and what constraints they impose on the remaining
candidates~\citep{Sandeman2026RAQM}.

% ─────────────────────────────────────────────────────────────────────────────
\section{The Experimental Programme}
\label{sec:experiment}
% ─────────────────────────────────────────────────────────────────────────────

Relational Actualism makes a large number of specific, quantitative,
falsifiable predictions. Table~\ref{tab:masses} collects the mass
predictions that follow from the cascade formula
$m_p = m_P\,\alpha_{\mathrm{EM}}^5/2^{28}$ (RASM Derivation~2).
Table~\ref{tab:predictions} lists further predictions.

\begin{table}[ht]
\centering
\caption{Mass and radius predictions from the $\mu^5$ cascade
(RASM Derivation~2). Inputs: $m_P$, $\alpha_{\mathrm{EM}}$ (IC30),
$L_q = 4$ (BDG, Lean-verified), $c_4 = 8$ (BDG).}
\label{tab:masses}
\renewcommand{\arraystretch}{1.3}
\begin{tabular}{lllrl}
\toprule
Quantity & Formula & Predicted & Observed & Error \\
\midrule
$m_p$ & $m_P\,\alpha_{\mathrm{EM}}^5/2^{28}$ & 941~MeV & 938~MeV & $0.3\%$ \\
$m_n$ & $= m_p$ (leading) & 941~MeV & 940~MeV & $0.2\%$ \\
$r_p$ & $L_q \times l_C$ & 0.841~fm & 0.841~fm & $0.03\%$ \\
$m_\pi$ & $(2/3)^5\,m_p$ & 124~MeV & 140~MeV & $11\%$ \\
$m_W$ & $6^{5/2}\,m_p$ & 83~GeV & 80~GeV & $3.3\%$ \\
$m_H$ & $(\alpha_{\mathrm{EM}}^{-1} - L_q)\,m_p$ & 125.2~GeV & 125.3~GeV & $0.06\%$ \\
\bottomrule
\end{tabular}
\end{table}

\begin{table}[ht]
\centering
\caption{Falsifiable predictions of Relational Actualism.}
\label{tab:predictions}
\renewcommand{\arraystretch}{1.3}
\begin{tabular}{p{5.2cm}p{4.5cm}p{3.5cm}}
\toprule
Prediction & Value / Status & Timescale \\
\midrule
BMV gravity-mediated entanglement & Strict null result & Near-term \\
No WIMP above neutrino floor; lensing from CGM halos & WIMP prohibited structurally; 4 lensing regimes & Ongoing \\
No fourth generation & $N_5$ absent in 4D BDG & LHC ongoing \\
No proton decay & LLC exact & Ongoing \\
Eridanus supervoid $H$ & $\approx 76.8~\mathrm{km\,s^{-1}\,Mpc^{-1}}$ & $\sim 5$ years \\
Hubble tension parameter-free & $H = H_{\mathrm{CMB}}(1+0.463|\delta|)$ & DESI \\
Neutrino mass sum & $\sum m_\nu \approx 59~\mathrm{meV}$ & CMB-S4/Euclid \\
$\sin^2\theta_W$ at GUT scale & $3/8$ & Existing data \\
$d = 4$ unique (four arguments) & Proposition & RASM \\
Stable matter only in $d=4$ & From $(4/3)d! = d\times 2^{d-1}$ & RASM \\
$\theta_{\mathrm{QCD}} = 0$ exactly; no axion & Proposition & RASM \\
$|V_{us}| = \sin(2/9) = 0.2204$ & Conjecture (1.8\% error) & RASM Conj.~7.3 \\
GW dispersion at Planck scale & Discrete corrections & Future GW detectors \\
CMB axis-of-evil persists in CMB-S4 & ${>}\,3\sigma$; structural, not statistical & CMB-S4 \\
CMB $(2/\ell)^2$ multipole hierarchy & Decreasing alignment for $\ell \geq 4$ & CMB-S4 \\
Parent SMBH mass & $\approx 2.5$--$6.6 \times 10^6\,M_\odot$ & RAGC \S6.3 \\
Near-minimum inflation & $N_{\mathrm{total}} \approx 64$--65 & Tensor-to-scalar ratio \\
Cosmic web--CMB axis correlation & Filaments aligned with spin axis & DESI/Euclid \\
No heat death & Fragmentation via dual severance & RAGC \S6.6 \\
Boltzmann Brains categorically impossible & Structural + dynamical & RAGC \S6.7 \\
$\alpha_{\mathrm{EM}}^{-1} = 137.036$ & Discrete Dyson equation, 0.00001\% & RASM \\
\bottomrule
\end{tabular}
\end{table}

The most decisive near-term test is the Bose-Marletto-Vedral
(BMV) protocol~\citep{Bose2017, Marletto2017}: two masses placed in
spatial superposition, entangled gravitationally. If gravity-mediated
entanglement is observed, the fundamental premise of RA --- that the
spacetime metric is updated strictly from actualized vertices --- is
falsified. RA predicts a strict null result.

The most important medium-term test is the neutrino mass sum.
Several forthcoming experiments (CMB-S4, Euclid, DESI) will reach
sensitivity $\sum m_\nu \sim 20~\mathrm{meV}$. RA predicts
$\sum m_\nu \approx 59~\mathrm{meV}$. If the bound falls below
this, the $SU(3)_{\mathrm{gen}}$ structure requires revision.

% ─────────────────────────────────────────────────────────────────────────────
\section{Open Derivations}
\label{sec:open}
% ─────────────────────────────────────────────────────────────────────────────

The programme originally identified four derivation targets. All four
have been resolved or dissolved. None remain open.

\paragraph{D1: Coupling constants --- resolved.}
The BDG integers $(1,-1,9,-16,8)$ are the \emph{unique} growth rule
for a self-consistent $d=4$ causal graph: the O14 uniqueness theorem
establishes that binomial inversion (M\"obius inversion on the Boolean
lattice) applied to the $d=4$ chord diagram moment sequence
$r_k = (2k+3)(2k+2)(2k+1)/6$~(Yeats 2025) yields exactly these
integers with zero free parameters (Rota 1964). The coupling constants
follow: $\alpha_{\mathrm{EM}}^{-1} = 137$ (Lean-verified, L08) with
dressed value $137.036$ via the discrete Dyson equation (IC30, 0.00001\%
of PDG); $\alpha_s(m_Z) = 1/\sqrt{72} = 0.11785$ (0.13\% of PDG).
No MCMC simulation is needed --- the coupling constants are uniquely
determined by the BDG integer structure.

\paragraph{D2: Gauge group structure --- resolved.}
The Standard Model gauge group $\mathrm{SU}(3) \times \mathrm{SU}(2)
\times \mathrm{U}(1)$ emerges from the alternating sign pattern of the
BDG coefficients $(+,-,+,-,+)$, which is forced by inclusion--exclusion
counting of causal intervals (GS02, confirmed by exact
enumeration~[CV]; see RATM~\S7.6). The $2+1$ split at depth~1
($E[N_1]=1$ at $\mu=1$, selecting one spatial direction) produces
$\mathrm{SU}(2)_L \times \mathrm{U}(1)_Y$; the depth-3 and depth-4
sectors produce $\mathrm{SU}(3)_c$. The remaining open step is
deriving the SU(3)$_{\mathrm{gen}}$ breaking pattern that explains why
Koide holds exactly for leptons but approximately for quarks.

\paragraph{D3: Lean formal proofs --- resolved.}
The BDG particle classification, colour confinement, and closure theorem
were already Lean-verified (\texttt{RA\_D1\_Proofs.lean}).
The chirality theorem (maximal parity violation from DAG acyclicity +
BDG filter) and baryon conservation (N$_2$ preserved in all asymptotic
extensions, verified exhaustively) are now proved in
\texttt{RA\_BaryonChirality.lean} (13 theorems, zero sorry tags,
compiled April 2026). The weak force is left-handed because time has a
direction: DAG acyclicity forbids backward-pointing edges, and the BDG
integers make the backward-winding pattern $N=(1,1,1,0)$ unstable
(score $= -7 < 0$).

\paragraph{D4: Type III$_1$ AQFT extension --- dissolved.}
The finite-dimensional Lean proofs of frame independence and Rindler
stationarity are the fundamental results in RA's discrete ontology.
The type~III$_1$ von Neumann algebra framework (Tomita-Takesaki modular
theory) is the continuum approximation to these discrete facts---it is
what the Lean-verified results \emph{look like} when coarse-grained
into continuum AQFT at $\mu \gg 1$. Showing compatibility is a
consistency check with the same dissolved status as O10 (Bianchi = LLC
in continuum costume) and O11 (Lorentz = causal invariance in continuum
costume). It is not a foundational requirement for the framework.

The history of the programme is that each target has yielded to an
RA-native insight once stated precisely in the framework's own terms.
D4 is the latest example: the ``problem'' was asking the discrete
theory to prove compatibility with a continuum approximation, which
inverts the ontological priority.

% ─────────────────────────────────────────────────────────────────────────────
\section{Further Consequences}
\label{sec:further_consequences}
% ─────────────────────────────────────────────────────────────────────────────

Four additional long-standing problems dissolve immediately once stated
in RA-native terms. Full derivations appear in RASM~\citep{Sandeman2026RASM}.

\paragraph{The strong CP problem: $\theta_{\mathrm{QCD}} = 0$ exactly.}
The QCD vacuum angle $\theta_{\mathrm{QCD}}$ is generated in the Standard
Model by instantons — field configurations that tunnel between topological
vacuum sectors. Instantons require closed loops in Euclidean spacetime.
In RA, time is the direction of causal edges; a closed Euclidean loop
requires a closed causal loop. The DAG is acyclic: instantons do not exist.
$\theta_{\mathrm{QCD}} = 0$ exactly, not merely small. A second independent
argument: the LLC on $N_2$ spatial charge is symmetric under charge conjugation
by construction, making CP an exact symmetry of the strong-force LLC.
The axion was introduced to relax $\theta_{\mathrm{QCD}}$ to zero dynamically.
RA predicts the axion does not exist — consistent with all observations.

\paragraph{The black hole information paradox: dissolved.}
A black hole is a Causal Severance: a topological partition of the causal
graph at which outgoing edges are absent. The LLC holds independently on
each side. Information is not lost; it is partitioned — encoded in the
initial LLC conditions of the daughter spacetime. An exterior observer
traces over the interior degrees of freedom and sees thermal (Hawking)
radiation, not because information is destroyed but because causal
disconnection produces a mixed state for any single-component observer.
The total system is unitary. The paradox arises only if one incorrectly
requires a single exterior observer to access causally disconnected
information — which Causal Severance structurally forbids.

\paragraph{The fine structure constant as self-consistency eigenvalue (conjecture).}
The per-link actualization probability satisfies $p_{\mathrm{th}} = \alpha_{\mathrm{EM}}$
via the optical theorem (proved, RASM): the fraction of EM interaction candidates
satisfying the actualization criterion $\Delta S > 0$ equals $\alpha_{\mathrm{EM}}$.
Combined with the score-$1$ floor theorem (RATM, Lean-proved), this gives
the self-consistency condition $\mu_{\mathrm{stable}} = 1/\alpha_{\mathrm{EM}} \approx 137$.

The empirical Wyler formula~\citep{Wyler1969} $\alpha = 9/(8\pi^4)(\pi^5/1920)^{1/4}$
(accurate to $0.00006\%$) decomposes with all factors BDG-derived: $|c_{N_2}| = 9$,
$d = 4$, and $\Gamma(5/4)$ from the causal diamond volume integral --- the same
$\Gamma(5/4)$ that appears in the BDG locality lemma (RACL). Two of the three
conditions needed to derive $\alpha$ from BDG integers alone are established;
the third (computing $V(SU(3)_{\mathrm{gen}})$ in BDG normalisation) is the
remaining open step (RASM Conjecture~7.2).

\paragraph{The Cabibbo angle from $SU(3)_{\mathrm{gen}}$ (conjecture).}
The Koide fit converges to $\theta_{\mathrm{lep}} = 2/9$ to one part in $10^5$,
where $2/9 = 2/(N_{\mathrm{gen}} \cdot N_{\mathrm{col}}) = 2/(3\times 3)$.
This predicts the Cabibbo angle: $|V_{us}| = \sin(2/9) = 0.2204$,
compared to the observed $0.2245$ ($1.8\%$ error).
The three quark Koide angles form an arithmetic progression
$\theta_0 \pm \delta$ around $\theta_0 = 2/9$,
with $\delta \approx 0.0665~\mathrm{rad}$ controlling the full CKM hierarchy.
Two numbers --- $2/9$ and $\delta$ --- parametrize all three CKM mixing angles.
(RASM Conjecture~7.3.)

\paragraph{Why $d = 4$: four independent uniqueness arguments.}
The BDG action in $d$ dimensions has $d/2 + 1$ non-gravitational degrees
of freedom. Argument~1: only $d = 4$ gives exactly three gauge forces.
Argument~2 (D4U02): the selectivity ceiling $\Delta S^*(\mu)$ has its
maximum at $\mu \approx 1$ only in $d = 4$.
Argument~3 (new): the cascade exponent $2N_c - 1$ equals the confinement
cycle length $d + 1$ only in $d = 4$.
Argument~4 (new): the Alexandrov/sphere geometric factor $(4/3)\,d!$
equals the BDG confinement parameter $d \times 2^{d-1}$ only in $d = 4$;
this is what allows hadronic self-consistency (RASM
Proposition~\ref*{prop:geometric_bdg}). Stable composite
matter --- protons, neutrons, atoms --- can only exist in $d = 4$.

\paragraph{The hierarchy problem: dissolved.}
The proton mass $m_p \approx 10^{-20}\,m_P$ is not ``unnaturally small.''
It is the \emph{unique} mass satisfying the BDG self-consistency condition
in $d = 4$: the cascade formula $m_p = m_P\,\alpha_{\mathrm{EM}}^5/2^{28}$
follows from the gluon vertex probability at the self-consistent internal
density $\mu_p = \alpha_{\mathrm{EM}}/32$, with $N_{\mathrm{eff}} = L_q^3$
derived from the $d = 4$ geometric identity. There is no fine-tuning because
there is no free parameter to tune; the BDG integers fix everything.

\paragraph{The Higgs mechanism in RA.}
In RA, the ``Higgs field'' is the background depth-2 actualization density.
The $W/Z$ bosons are marginal BDG vertices with score $S(1,0,0,0) = c_0 + c_1
= 0$ exactly; they require depth-2 dressing to actualize. The photon escapes
as $N = (1,1,0,0)$ with $S = 9$, actualizing freely. The Higgs boson is a
coherent perturbation of 133 background modes, with mass $m_H = 133\,m_p
= 125.2~\mathrm{GeV}$ (RASM equation~\ref*{eq:higgs_mass}).
No fundamental scalar field is needed.

\paragraph{The baryon asymmetry as a Causal Severance initial condition.}
Baryon number is exactly conserved in RA (Proposition~5.1 in RASM).
Sakharov's first condition — baryon number violation — is therefore
absent. This does not prevent RA from accounting for $\eta$; it means
dynamic baryogenesis does not occur. The Big Bang is a Causal Severance
nucleation: our universe inherited its net baryon number $B \neq 0$
from the LLC boundary conditions at that event. The baryon asymmetry
$\eta \approx 6.1\times 10^{-10}$~\citep{PDG2022} is an initial
condition set at the Planck scale, not a dynamic outcome. The prediction:
no signatures of electroweak baryogenesis in CMB data; the
baryon-to-photon ratio is isotropic at all scales.

% ─────────────────────────────────────────────────────────────────────────────
\section{Lean 4 Formal Verification}
\label{sec:lean}
% ─────────────────────────────────────────────────────────────────────────────

The RA programme uses Lean~4/Mathlib as a proof assistant to machine-verify
the discrete graph-theoretic core of the framework. This is not cosmetic:
the Lean proofs establish that the claims about the growing DAG are
internally consistent at the level of formal mathematics, independently
of any physical interpretation.

The current Lean corpus comprises \textbf{176 verified theorems across seven
proof files}, with one intentional sorry (the LQI adapter in
\texttt{RA\_AQFT\_Proofs\_v10.lean}) and no axioms beyond Mathlib:

\begin{center}
\renewcommand{\arraystretch}{1.2}
\begin{tabular}{p{5.5cm}p{4cm}p{2.5cm}}
\toprule
Result & File & Status \\
\midrule
Local Ledger Condition (L01) & \texttt{RA\_GraphCore.lean} & Proved, 0 sorry \\
Graph Cut Theorem (L02) & \texttt{RA\_GraphCore.lean} & Proved, 0 sorry \\
Markov Blanket Shielding (L03) & \texttt{RA\_GraphCore.lean} & Proved, 0 sorry \\
Amplitude locality (O01) & \texttt{RA\_AmpLocality.lean} & Proved, 0 sorry \\
Causal invariance (O02) & \texttt{RA\_AmpLocality.lean} & Proved, 0 sorry \\
Frame independence of $\Delta S(\rho\|\sigma_0)$ & \texttt{RA\_AQFT\_Proofs\_v10.lean} & Proved, 0 sorry \\
Rindler stationarity (Unruh resolution) & \texttt{RA\_AQFT\_Proofs\_v10.lean} & Proved, 0 sorry \\
Koide $K = 2/3$ identity (L09) & \texttt{RA\_Koide.lean} & Proved, 0 sorry \\
$SU(3)_{\mathrm{gen}}$ coherent state theorem & \texttt{RA\_D1\_Proofs.lean} & Proved, 0 sorry \\
BDG particle classification (5 types, L11) & \texttt{RA\_D1\_Proofs.lean} & Proved, 0 sorry \\
BDG closure theorem (124 cases) & \texttt{RA\_D1\_Proofs.lean} & Proved, 0 sorry \\
Colour confinement ($L{=}3$, $L{=}4$; L10) & \texttt{RA\_D1\_Proofs.lean} & Proved, 0 sorry \\
$\alpha_{\mathrm{EM}}^{-1} = 137$ (L08) & \texttt{RA\_D1\_Proofs.lean} & Proved, 0 sorry \\
$\Pact$ conservation theorem & \texttt{RA\_D1\_Proofs.lean} & Structural proof \\
BDG locality lemma & numerically verified & Confirmed \\
\bottomrule
\end{tabular}
\end{center}

\noindent
Amplitude locality (O01) is proved as a theorem of discrete DAG
dynamics---the BDG action increment depends only on the causal past
by transitivity of the causal order---making causal invariance (O02)
unconditional with zero sorry tags. The type~III$_1$ AQFT framework
(Tomita-Takesaki modular theory, von Neumann algebras) is the continuum
approximation to these discrete results; in RA's ontology, the
finite-dimensional Lean proofs are fundamental and compatibility with
the continuum is a consistency check, not a foundational requirement
(same dissolved status as O10 and O11).

A Python notebook independently confirms 52/52 numerical checks at
machine precision, covering the Koide formula, the BDG coefficient
ratios, the Wyler formula decomposition, and the coupling constant
hierarchies. The machine verification strategy --- using Lean for
structural results and Python for numerical confirmation --- provides
two independent lines of evidence for the quantitative claims of the suite.

% ─────────────────────────────────────────────────────────────────────────────
\section{Prior Approaches and the Position of RA}
\label{sec:prior}
% ─────────────────────────────────────────────────────────────────────────────

\paragraph{General relativity and quantum field theory.}
Both frameworks are derived in RA as macroscopic approximations to the
causal graph dynamics. They are not competing descriptions of nature;
they are Level~3 descriptions that apply in different density regimes
of the same Level~1 algorithm.

\paragraph{Loop Quantum Gravity.} LQG derives geometry from a spin
network background but imports the Standard Model as separate matter
content. It offers no account of why the gauge group has three factors
or why charge is quantised. RA requires no background; the Standard
Model structure is a consequence of 4D BDG topology.

\paragraph{String Theory.} String theory can reproduce a rich particle
spectrum but requires 10 or 11 spacetime dimensions and a landscape of
$\sim 10^{500}$ vacua with no selection principle. RA is uniquely
4-dimensional by BDG + Lovelock, has no landscape, and no free
parameters at the topological level.

\paragraph{Causal set theory.} RA is built directly on the causal set
foundation of Sorkin~\citep{Sorkin1994} and Bombelli~\citep{Bombelli1987},
extending it with three new elements: (1) the BDG actualization filter
$\SBDG(v) > 0$ as a physical selection criterion (rather than a dynamics
to be inferred from path integrals); (2) the LLC as a vertex-level
conservation law, Lean-verified; (3) the Poisson-CSG growth model as
the mean-field description of the stochastic graph rewriting.
The key advance over standard causal set theory is the BDG locality lemma
(RACL): the dense-graph limit of the BDG sum is provably local, which
is what allows Lovelock's uniqueness argument to apply. Without locality,
Lovelock does not constrain the field equation; with it, Einstein is forced.
Causal set theory asks: can GR emerge from a causal set? RA answers: it
must, and here is a machine-checkable proof.

The Rideout-Sorkin Classical Sequential Growth model~\citep{Rideout2000}
is the direct precursor to the RA Poisson-CSG: the CSG dynamics supply
the Born-rule-like probability measure for graph growth that RA uses.
RA extends CSG by adding the BDG filter and identifying the LLC
as the conservation law that forces GR.

\paragraph{Bilson-Thompson topological preons.}
Bilson-Thompson's braid model~\citep{BilsonThompson2005} maps the
first-generation Standard Model onto ribbon braids. RA is in the same
direction but derives the stability selection rule (the BDG filter)
that the braid model lacks. Whether RA's stable patterns correspond to
Bilson-Thompson braids is an open question that D1 would resolve.

% ─────────────────────────────────────────────────────────────────────────────
\section{Conclusion}
\label{sec:conclusion}
% ─────────────────────────────────────────────────────────────────────────────

Physics has spent a century seeking the reconciliation of quantum
mechanics and general relativity, and asking why the Standard Model
has the structure it has. Relational Actualism proposes that both
questions are symptoms of a shared error: treating spacetime,
matter, and quantum states as primitive, when none of them is.

The primitive is the actualization event. Not the particle, not
the field, not the wave function, not the spacetime manifold.
The moment at which a physical interaction becomes irreversible ---
at which $\Delta S(\rho\|\sigma_0) > 0$ and a permanent causal edge
is written into the DAG.

From this:

\medskip
\begin{center}
\renewcommand{\arraystretch}{1.2}
\begin{tabular}{ll}
\toprule
Result & Status \\
\midrule
$c = l_P/t_P$; $E = \gamma mc^2$; $\tau = N_{\mathrm{act}}\cdot t_P$ & Proved \\
GR inevitable (BDG uniqueness + Lovelock) & Proved \\
Three non-gravitational forces in 4D & Proved \\
Charge quantisation in units of $e/3$ & Proved \\
Photon masslessness & Proved \\
Colour confinement & Proved \\
Exact baryon conservation & Proved (Lean-backed) \\
Maximal parity violation & Proved \\
Exactly three generations & Proved \\
Koide lepton mass formula & Conjectured; $K = 2/3$ exact \\
GUT unification at Planck scale & Conjecture \\
Flat rotation curves; Hubble tension & Derived \\
$\kappa(\theta)$ Bullet Cluster profile & Qualitative (RADM §4.2) \\
4-regime lensing classification & Established (RADM §4.3) \\
DJW lensing $\kappa(b) = v_\mathrm{flat}^2/(4G\Sigma_\mathrm{cr}b)$ & Derived (RADM §4.1) \\
$\kappa(b)$ cutoff at $R_\mathrm{baryon}$ vs NFW & Falsifiable prediction \\
Galaxy-galaxy amplitude (CGM) & Prediction, open calc \\
Origin of life (Causal Firewall) & Derived \\
QFT fault-tolerance threshold & Derived \\
Measurement problem & Dissolved \\
$\theta_{\mathrm{QCD}} = 0$ exactly; no axion & Proved \\
Black hole information paradox & Dissolved \\
Baryon asymmetry as Causal Severance IC & Derived \\
$d = 4$ spacetime dimensionality & Conjecture \\
Quark Koide breaking (RA form; $\alpha_s$ external input) & Conjecture \\
Wyler formula for $\alpha_{\mathrm{EM}}$ (0.00006\%) & Conjecture \\
$\theta_{\mathrm{lep}} = 2/9$ (Koide fit, $1:10^5$) & Conjecture \\
$|V_{us}| = \sin(2/9)$ Cabibbo prediction (1.8\%) & Conjecture \\
$\delta = \tfrac{3}{4}\alpha_s/\pi$ ($0.1\%$) & Proved \\
$|V_{cb}| = (2/\pi)\delta$ ($0.4\%$) & Conjecture \\
\bottomrule
\end{tabular}
\end{center}

\medskip
The question driving the next phase of the programme is exact:
is the complete set of $\SBDG > 0$-stable self-similar patterns
in a 4D BDG-filtered Poisson-CSG precisely the first-generation
Standard Model particle spectrum? If yes --- if the BDG-filter MCMC
confirms that the stable patterns are exactly $\{e, u, d, \nu_e,
\gamma, g, W, Z, H\}$ --- then all of physics follows from the
single fact that some things happen and others do not.

\medskip
The programme as it stands in early 2026:
\begin{itemize}[noitemsep]
  \item RAQM at \emph{Foundations of Physics} (awaiting editor assignment).
  RACL and RATM targeting \emph{Physical Review D} (revised after initial rejections).
  RACF at \emph{International Journal of Astrobiology} (under review).
  RACL targeting \emph{Physical Review D} (revised after CQG desk rejection).
  \item Twelve papers in the suite; eleven on Zenodo with permanent
  DOIs (RACF submitted to \emph{International Journal of Astrobiology}).
  \item 176 Lean-verified theorems across eight proof files, one
  intentional sorry (LQI adapter).
  \item 52/52 independent numerical checks at machine precision.
  \item All four derivation targets resolved or dissolved:
  D1 (coupling constants), D2 (gauge groups), D3 (chirality/baryon,
  compiled April 2026), D4 (type~III$_1$ AQFT, dissolved).
\end{itemize}

The programme does not claim completeness. It claims internal consistency
at the discrete level (machine-verified), correct macroscopic limits
(GR inevitable, Standard Model derived), and specific falsifiable
predictions (BMV null result; no WIMP; $\sum m_\nu \approx 59$~meV;
$H_{\mathrm{Eridanus}} \approx 76.8$~km/s/Mpc; $\kappa(b)$ cutoff at
$R_{\mathrm{baryon}}$). The framework is falsifiable. The search continues.

% ─────────────────────────────────────────────────────────────────────────────

\section*{Declarations}
\textbf{Funding:} Not applicable.\quad
\textbf{Competing interests:} The author declares no competing interests.

\textbf{Ethics statement:} This article does not contain any studies with human participants or animals performed by the author.

\textbf{AI use disclosure:} Claude (Anthropic, model \texttt{claude-sonnet-4-6}) was used in the preparation of this manuscript, including drafting and revision of text, mathematical derivation assistance, \LaTeX{} compilation, and Lean~4 proof development. The author is solely responsible for all scientific content, mathematical claims, and conclusions. Gemini (Google DeepMind, Gemini~2.0~Pro) was used for supplementary literature checking and cross-verification of symbolic calculations.

\textbf{Data availability:} All Lean~4 proof files and Python verification scripts are publicly available at \href{https://github.com/jsandeman/Relational-Actualism}{github.com/jsandeman/Relational-Actualism}.

\textbf{Preprint availability:} \href{https://doi.org/10.5281/zenodo.19229438}{doi.org/10.5281/zenodo.19229438}.

\bibliographystyle{plainnat}

\begin{thebibliography}{99}

\bibitem[Sandeman(2026RACL)]{Sandeman2026RACL}
J.~F. Sandeman.
\newblock {Relational Actualism: The Einstein Field Equations as the Unique
  Macroscopic Limit of the Actualization Graph (RACL)}.
\newblock \textit{Preprint}, 2026.
\newblock \href{https://doi.org/10.5281/zenodo.19270020}{doi.org/10.5281/zenodo.19270020}.


\bibitem[Sandeman(2026RATM)]{Sandeman2026RATM}
J.~F. Sandeman.
\newblock {Relational Actualism: The Topology of Matter (RATM)}.
\newblock \textit{Preprint}, 2026.
\newblock \href{https://doi.org/10.5281/zenodo.19265651}{doi.org/10.5281/zenodo.19265651}.

\bibitem[Sandeman(2026RAEB)]{Sandeman2026RAEB}
J.~F. Sandeman.
\newblock {Relational Actualism: The Engine of Becoming (RAEB)}.
\newblock \textit{Preprint}, 2026.
\newblock \href{https://doi.org/10.5281/zenodo.19198120}{doi.org/10.5281/zenodo.19198120}.


\bibitem[Benincasa \& Dowker(2010)]{Benincasa2010}
D.~M.~T.~Benincasa and F.~Dowker.
\newblock The scalar curvature of a causal set.
\newblock \emph{Physical Review Letters}, 104(18):181301, 2010.
\newblock \href{https://doi.org/10.1103/PhysRevLett.104.181301}%
  {doi:10.1103/PhysRevLett.104.181301}.

\bibitem[Bilson-Thompson(2005)]{BilsonThompson2005}
S.~O.~Bilson-Thompson.
\newblock A topological model of composite preons.
\newblock \emph{arXiv preprint}, 2005.
\newblock \href{https://arxiv.org/abs/hep-ph/0503213}{arXiv:hep-ph/0503213}.

\bibitem[Bombelli et~al.(1987)]{Bombelli1987}
L.~Bombelli, J.~Lee, D.~Meyer, and R.~D.~Sorkin.
\newblock Space-time as a causal set.
\newblock \emph{Physical Review Letters}, 59(5):521--524, 1987.
\newblock \href{https://doi.org/10.1103/PhysRevLett.59.521}%
  {doi:10.1103/PhysRevLett.59.521}.

\bibitem[Bose et~al.(2017)]{Bose2017}
S.~Bose et~al.
\newblock Spin entanglement witness for quantum gravity.
\newblock \emph{Physical Review Letters}, 119(24):240401, 2017.
\newblock \href{https://doi.org/10.1103/PhysRevLett.119.240401}%
  {doi:10.1103/PhysRevLett.119.240401}.

\bibitem[Koide(1983)]{Koide1983}
Y.~Koide.
\newblock New view of quark and lepton mass hierarchy.
\newblock \emph{Physical Review D}, 28(1):252--254, 1983.
\newblock \href{https://doi.org/10.1103/PhysRevD.28.252}%
  {doi:10.1103/PhysRevD.28.252}.

\bibitem[Lovelock(1971)]{Lovelock1971}
D.~Lovelock.
\newblock The Einstein tensor and its generalizations.
\newblock \emph{Journal of Mathematical Physics}, 12(3):498--501, 1971.
\newblock \href{https://doi.org/10.1063/1.1665613}{doi:10.1063/1.1665613}.

\bibitem[Marletto \& Vedral(2017)]{Marletto2017}
C.~Marletto and V.~Vedral.
\newblock Gravitationally induced entanglement between two massive particles
  is sufficient evidence of quantum effects in gravity.
\newblock \emph{Physical Review Letters}, 119(24):240402, 2017.
\newblock \href{https://doi.org/10.1103/PhysRevLett.119.240402}%
  {doi:10.1103/PhysRevLett.119.240402}.

\bibitem[Particle Data Group(2022)]{PDG2022}
R.~L.~Workman et~al.
\newblock Review of particle physics.
\newblock \emph{Progress of Theoretical and Experimental Physics},
  2022(8):083C01, 2022.
\newblock \href{https://doi.org/10.1093/ptep/ptac097}{doi:10.1093/ptep/ptac097}.

\bibitem[Sandeman(2026a)]{Sandeman2026RAQM}
J.~F.~Sandeman.
\newblock Relational Actualism and Quantum Mechanics (RAQM).
\newblock \emph{Foundations of Physics} (under review), 2026.
\newblock \href{https://doi.org/10.5281/zenodo.19174380}%
  {doi:10.5281/zenodo.19174380}.

\bibitem[Sandeman(2026b)]{Sandeman2026RAGC}
J.~F.~Sandeman.
\newblock Relational Actualism: Gravity and Cosmology (RAGC).
\newblock Preprint, 2026.
\newblock \href{https://doi.org/10.5281/zenodo.19197999}%
  {doi:10.5281/zenodo.19197999}.

\bibitem[Sandeman(2026d)]{Sandeman2026RAHC}
J.~F.~Sandeman.
\newblock Relational Actualism: Algorithmic Causal Dynamics (RAHC).
\newblock Preprint, 2026.
\newblock \href{https://doi.org/10.5281/zenodo.19198171}%
  {doi:10.5281/zenodo.19198171}.

\bibitem[Sandeman(2026e)]{Sandeman2026RACI}
J.~F.~Sandeman.
\newblock Relational Actualism: Complexity and Intelligence (RACI).
\newblock Preprint, 2026.
\newblock \href{https://doi.org/10.5281/zenodo.19198224}%
  {doi:10.5281/zenodo.19198224}.

\bibitem[Sandeman(2026e2)]{Sandeman2026RACF}
J.~F.~Sandeman.
\newblock The Causal Firewall: Substrate-Independent Physical Constraints
  on the Distribution of Life in the Universe (RACF).
\newblock \textit{Int.\ J.\ Astrobiol.} (submitted), 2026.

ewblock \href{https://doi.org/10.5281/zenodo.19319374}{doi:10.5281/zenodo.19319374}.

\bibitem[Sandeman(2026f)]{Sandeman2026RAQI}
J.~F.~Sandeman.
\newblock Relational Actualism: Implications for Quantum Information (RAQI).
\newblock Preprint, 2026.
\newblock \href{https://doi.org/10.5281/zenodo.19198285}%
  {doi:10.5281/zenodo.19198285}.

\bibitem[Sandeman(2026g)]{Sandeman2026RADM}
J.~F.~Sandeman.
\newblock Relational Actualism: Dark Matter as Actualization-Density
  Topology (RADM).
\newblock Preprint, 2026.
\newblock \href{https://doi.org/10.5281/zenodo.19198513}%
  {doi:10.5281/zenodo.19198513}.

\bibitem[Sandeman(2026h)]{Sandeman2026RASM}
J.~F.~Sandeman.
\newblock Relational Actualism and the Standard Model (RASM).
\newblock Preprint, 2026.
\newblock \href{https://doi.org/10.5281/zenodo.19229335}{doi.org/10.5281/zenodo.19229335}

\bibitem[Sorkin(1994)]{Sorkin1994}
R.~D.~Sorkin.
\newblock Quantum mechanics as quantum measure theory.
\newblock \emph{Modern Physics Letters A}, 9(33):3119--3127, 1994.
\newblock \href{https://doi.org/10.1142/S021773239400294X}%
  {doi:10.1142/S021773239400294X}.

\bibitem[Rideout \& Sorkin(2000)]{Rideout2000}
D.~P.~Rideout and R.~D.~Sorkin.
\newblock A classical sequential growth dynamics for causal sets.
\newblock \emph{Physical Review D}, 61(2):024002, 2000.
\newblock \href{https://doi.org/10.1103/PhysRevD.61.024002}%
  {doi:10.1103/PhysRevD.61.024002}.

\bibitem[Wyler(1969)]{Wyler1969}
A.~Wyler.
\newblock On the conformal groups in the theory of relativity and their
  unitary representations.
\newblock \emph{Comptes Rendus Acad.~Sci.~Paris}, 269:743--745, 1969.

\bibitem[Walker et~al.(2023)]{Walker2023}
S.~I.~Walker et~al.
\newblock Explaining the origin of life is harder than we think.
\newblock \emph{arXiv preprint}, 2023.
\newblock \href{https://arxiv.org/abs/2308.09738}{arXiv:2308.09738}.

\end{thebibliography}

\end{document}
